% Источники: raw/rk2-additions/photo_2026-11-07_23-15-50_1002834.jpg
% -- photo_2026-11-07_23-15-51_1002843.jpg, по возрастанию номера.
% Компиляция: tectonic rk2-additions.tex
% Сокращения раскрыты. Существенные ошибки исходника отмечены явно.
\ifdefined\lessonpreambleloaded
  \expandafter\endinput
\fi
\let\lessonpreambleloaded\empty
\documentclass[a4paper,11pt]{article}
\usepackage{fontspec}
\setmainfont{Times New Roman}
\usepackage[russian]{babel}
\usepackage{amsmath,amssymb}
\usepackage[margin=22mm]{geometry}
\usepackage{indentfirst}
\usepackage{needspace}
\usepackage{tikz}
\usetikzlibrary{patterns}
\usepackage[hidelinks]{hyperref}
\usepackage{fancyhdr}
\pagestyle{fancy}
\fancyhf{}
\renewcommand{\headrulewidth}{0pt}
\fancyfoot[C]{\footnotesize\url{https://github.com/bmstu-software-engineering/cs-s5-tv-26y}}
\fancyfoot[R]{\thepage}
\setlength{\emergencystretch}{2em}
\newcommand{\sourcepage}[1]{\par\Needspace{6\baselineskip}\medskip\noindent{\footnotesize\itshape Источник: \detokenize{#1}.}\par\nopagebreak\smallskip}
\newcommand{\editorial}[1]{\par\smallskip\noindent{\small\itshape Примечание к переносу. #1}\par\smallskip}
\newcommand{\lessontitle}[3]{\begin{center}{\large\bfseries Теория вероятностей. #1}\\[4pt]{\large #2}\\[4pt]#3\end{center}}
\newcommand{\parttitle}[1]{\par\Needspace{4\baselineskip}\medskip\noindent\textbf{#1}\par\nopagebreak\smallskip}
\tikzset{element/.style={draw,rectangle,minimum width=7mm,minimum height=5mm,fill=white},every picture/.style={line width=.5pt}}

\begin{document}
\lessontitle{Дополнительные материалы к РК2}{Комбинаторика, надёжность, полная вероятность и формула Байеса}{}

\sourcepage{rk2-additions/photo_2026-11-07_23-15-50_1002834.jpg}
\section*{Схема Бернулли. Исправные детали}
Проводятся $n=100$ независимых испытаний.
$A$ --- деталь бракованная (успех), $\overline A$ --- деталь исправна (неудача):
\[
 p=P(A)=0{,}015,\qquad q=P(\overline A)=1-p=0{,}985.
\]
$B$ --- все детали исправны. По формуле Бернулли
\[
 P(B)=P_{100}(0)=C_{100}^0p^0q^{100}=0{,}985^{100}.
\]
По формуле Пуассона при $\lambda=np=0{,}015\cdot100=1{,}5$
\[
 P(B)\approx\frac{1{,}5^0e^{-1{,}5}}{0!}=e^{-1{,}5}.
\]
\parttitle{Номера на полях}
В рамке записаны: 3.34, 3.31, 3.39, 3.40, 3.42.

\section*{Дополнительная пара. Шары в урне}
В урне $6$ чёрных и $14$ белых шаров, всего $20$. Вынимают $5$ шаров.
Чёрные шары пронумеруем от $1$ до $6$, белые --- от $7$ до $20$.

\sourcepage{rk2-additions/photo_2026-11-07_23-15-50_1002835.jpg}
Рассматриваются события:
\begin{itemize}
 \item $A$ --- вынуто ровно три чёрных шара;
 \item $B$ --- вынут хотя бы один чёрный шар.
\end{itemize}
Исход --- набор $\{x_1,x_2,x_3,x_4,x_5\}$ различных номеров шаров.
Порядок не важен, повторений нет:
\[
 \Omega=\bigl\{\{x_1,\ldots,x_5\}:x_i\in\{1,\ldots,20\},\
 x_i\ne x_j\text{ при }i\ne j\bigr\}.
\]
Это сочетания без повторений из $20$ по $5$:
\[
 N=|\Omega|=C_{20}^5=\frac{20!}{5!15!}.
\]
\parttitle{а) Ровно три чёрных шара}
В наборе можно обозначить номера чёрных шаров через $x_1,x_2,x_3$,
а белых --- через $x_4,x_5$:
\[
 x_i\in\{1,\ldots,6\}\quad(i=1,2,3),\qquad
 x_j\in\{7,\ldots,20\}\quad(j=4,5),
\]
причём все пять номеров различны.

\sourcepage{rk2-additions/photo_2026-11-07_23-15-50_1002836.jpg}
По основной формуле комбинаторики
\[
 N_A=|A|=C_6^3C_{14}^2=\frac{6!}{3!3!}\frac{14!}{2!12!},
\]
\[
 P(A)=\frac{N_A}{N}=\frac{C_6^3C_{14}^2}{C_{20}^5}
 =\frac{6!}{3!3!}\frac{14!}{2!12!}\frac{5!15!}{20!}
 =\frac{455}{3876}.
\]
\editorial{В последнем сокращении рукописи потерян множитель $6$:
записано $13\cdot7\cdot5/(4\cdot17\cdot19\cdot18)$.
Здесь итоговая дробь пересчитана по исходной формуле сочетаний.}

\parttitle{б) Хотя бы один чёрный шар}
Противоположное событие $\overline B$ --- ни одного чёрного шара:
\[
 \overline B=\bigl\{\{x_1,\ldots,x_5\}:x_i\in\{7,\ldots,20\},\
 x_i\ne x_j\text{ при }i\ne j\bigr\},
\]
\[
 N_{\overline B}=C_{14}^5=\frac{14!}{5!9!}.
\]
\[
 P(B)=1-P(\overline B)=1-\frac{N_{\overline B}}N
 =1-\frac{C_{14}^5}{C_{20}^5}
 =1-\frac{14!15!}{9!20!}\approx0{,}871.
\]
\editorial{В рукописи итоговое приближение указано как $0{,}775$.
Оно исправлено по записанным числам шаров и формуле.}

\section*{Четыре пассажира и девять вагонов}
Четыре пассажира садятся в поезд из девяти вагонов. Требуется найти вероятности:
\begin{enumerate}
 \item все сели в один чётный вагон;
 \item все сели в разные вагоны;
 \item по два человека сели в первый и второй вагоны.
\end{enumerate}
В заметке перед задачей подчёркнуто: в задачах о костях, лифтах и вагонах
важно учитывать порядок.

\sourcepage{rk2-additions/photo_2026-11-07_23-15-50_1002837.jpg}
Исход $(x_1,x_2,x_3,x_4)$, где $x_i\in\{1,\ldots,9\}$ --- номер вагона,
в который сел $i$-й человек. Порядок важен, повторения возможны:
\[
 \Omega=\{(x_1,x_2,x_3,x_4):x_i\in\{1,\ldots,9\},\ i=1,\ldots,4\}.
\]
Это размещения с повторениями из $9$ по $4$:
\[
 N=|\Omega|=\widetilde A_9^4=9^4.
\]
\parttitle{1. Один чётный вагон}
\[
 A=\{(x,x,x,x):x\in\{2,4,6,8\}\},\qquad N_A=4,
\]
\[
 P(A)=\frac{N_A}{N}=\frac4{9^4}.
\]
\parttitle{2. Разные вагоны}
\[
 B=\{(x_1,x_2,x_3,x_4):x_i\in\{1,\ldots,9\},\
 x_i\ne x_j\text{ при }i\ne j\}.
\]

\sourcepage{rk2-additions/photo_2026-11-07_23-15-50_1002838.jpg}
Порядок важен, повторений нет, поэтому
\[
 N_B=A_9^4=\frac{9!}{5!}=6\cdot7\cdot8\cdot9,\qquad
 P(B)=\frac{6\cdot7\cdot8\cdot9}{9^4}.
\]
\parttitle{3. По два человека в первом и втором вагонах}
Не важно, какие именно люди сели в каждый из этих вагонов:
\[
 C=\{(x_1,x_2,x_3,x_4):x_i\in\{1,2\},\
 |\{i:x_i=1\}|=2\}.
\]
Используем полиномиальный коэффициент:
\[
 N_C=C(2,2,0,0,0,0,0,0,0)=\frac{4!}{2!2!}
 =\frac{4\cdot3}{2}=6,
\]
\[
 P(C)=\frac{N_C}{N}=\frac6{9^4}.
\]

\section*{Задача 2. Соединения элементов}
$A$ --- система работает, $\overline A$ --- отказ системы;
$A_i$ --- работает $i$-й элемент, $\overline A_i$ --- его отказ.

\parttitle{Последовательное соединение}
\begin{center}
\begin{tikzpicture}
 \draw (0,0)--(5,0);
 \node[element] at (1,0){1};\node[element] at (2.5,0){2};\node[element] at (4,0){3};
 \fill (0,0) circle (1.3pt);\fill (5,0) circle (1.3pt);
\end{tikzpicture}
\end{center}
\[
 A=A_1A_2A_3=A_1\cap A_2\cap A_3.
\]

\sourcepage{rk2-additions/photo_2026-11-07_23-15-50_1002839.jpg}
\[
 \overline A=\overline A_1\cup\overline A_2\cup\overline A_3.
\]
При независимых отказах
\[
 P(A)=P(A_1)P(A_2)P(A_3),\qquad P(\overline A)=1-P(A).
\]
\parttitle{Параллельное соединение}
\begin{center}
\begin{tikzpicture}
 \draw (0,0)--(1,0) (1,-.8)--(1,.8) (3,-.8)--(3,.8) (3,0)--(4,0);
 \foreach \y/\n in {.8/1,0/2,-.8/3}{\draw (1,\y)--(3,\y);\node[element] at (2,\y){\n};}
 \fill (0,0) circle (1.3pt);\fill (4,0) circle (1.3pt);
\end{tikzpicture}
\end{center}
\[
 A=A_1\cup A_2\cup A_3,\qquad
 \overline A=\overline A_1\cap\overline A_2\cap\overline A_3,
\]
\[
 P(\overline A)=P(\overline A_1)P(\overline A_2)P(\overline A_3),\qquad
 P(A)=1-P(\overline A).
\]

\parttitle{Система из четырёх элементов}
Верхняя ветвь содержит элемент 1. Нижняя --- элемент 2 последовательно
с параллельным соединением элементов 3 и 4.
\begin{center}
\begin{tikzpicture}
 \draw (0,0)--(1,0) (1,-.8)--(1,.9) (1,.9)--(5,.9)
 (1,-.8)--(3,-.8) (3,-1.2)--(3,-.4)
 (3,-.4)--(5,-.4) (3,-1.2)--(5,-1.2)
 (5,-1.2)--(5,.9) (5,0)--(6,0);
 \node[element] at (3,.9){1};\node[element] at (2,-.8){2};
 \node[element] at (4,-.4){3};\node[element] at (4,-1.2){4};
\end{tikzpicture}
\end{center}
Вероятности работы:
\[
 P(A_1)=0{,}7,\quad P(A_2)=0{,}8,\quad
 P(A_3)=0{,}6,\quad P(A_4)=0{,}9.
\]
Вероятности отказа:
\[
 P(\overline A_1)=0{,}3,\quad P(\overline A_2)=0{,}2,\quad
 P(\overline A_3)=0{,}4,\quad P(\overline A_4)=0{,}1.
\]
\editorial{Рядом с вероятностями отказа в рукописи записано
$P(\overline A_i)=P(A_i)-1$. Знак исправлен:
$P(\overline A_i)=1-P(A_i)$; приведённые числовые значения соответствуют этой формуле.}

\sourcepage{rk2-additions/photo_2026-11-07_23-15-50_1002840.jpg}
Событие работы системы:
\[
 A=A_1\cup\bigl(A_2\cap(A_3\cup A_4)\bigr).
\]
Отказ системы:
\[
 \overline A=\overline A_1\cap
 \bigl(\overline A_2\cup(\overline A_3\cap\overline A_4)\bigr).
\]
При независимых отказах
\begin{align*}
 P(\overline A)
 &=P(\overline A_1)P\bigl(\overline A_2\cup(\overline A_3\cap\overline A_4)\bigr)\\
 &=P(\overline A_1)\bigl[1-P(A_2)P(A_3\cup A_4)\bigr]\\
 &=P(\overline A_1)\bigl[1-P(A_2)
 (1-P(\overline A_3)P(\overline A_4))\bigr]\\
 &=0{,}3\bigl[1-0{,}8(1-0{,}4\cdot0{,}1)\bigr]
 =0{,}0696\approx0{,}07.
\end{align*}
\[
 P(A)=1-P(\overline A)=0{,}9304\approx0{,}93.
\]
\editorial{В одной промежуточной строке исходника пропущены скобки
в выражении для отказа. Здесь они восстановлены по схеме и последующему вычислению.
Точные значения $0{,}0696$ и $0{,}9304$ поясняют округлённые ответы рукописи.}

\sourcepage{rk2-additions/photo_2026-11-07_23-15-50_1002841.jpg}
\section*{Задача 3. Экзаменационные билеты}
Двадцать экзаменационных билетов содержат по два вопроса.
Студент знает $30$ вопросов из $40$. Для сдачи экзамена нужно ответить
на оба вопроса билета либо на один вопрос билета и один дополнительный.
\begin{enumerate}
 \item Найти вероятность сдачи экзамена.
 \item Известно, что студент сдал экзамен. Найти вероятность того,
 что в своём билете он знал ровно один вопрос.
\end{enumerate}
\Needspace{8\baselineskip}
$A$ --- студент сдал экзамен. Выберем гипотезы:
\begin{itemize}
 \item $H_1$ --- не ответил ни на один вопрос своего билета;
 \item $H_2$ --- ответил ровно на один вопрос своего билета;
 \item $H_3$ --- ответил на оба вопроса своего билета.
\end{itemize}
Гипотезы попарно несовместны и образуют полную группу.
$B$ --- ответил на первый вопрос своего билета,
$C$ --- ответил на второй. Тогда $H_1=\overline B\cap\overline C$.

\sourcepage{rk2-additions/photo_2026-11-07_23-15-50_1002842.jpg}
\parttitle{Вероятности гипотез}
\[
 P(H_1)=P(\overline B)P(\overline C\mid\overline B)
 =\frac{10}{40}\frac9{39}=\frac{90}{40\cdot39}.
\]
\[
 H_2=(B\cap\overline C)\cup(\overline B\cap C).
\]
Эти два события несовместны, а ответы на вопросы при выборе без возвращения
зависимы. По формуле умножения
\begin{align*}
 P(H_2)&=P(B)P(\overline C\mid B)+P(\overline B)P(C\mid\overline B)\\
 &=\frac{30}{40}\frac{10}{39}+\frac{10}{40}\frac{30}{39}
 =\frac{600}{40\cdot39}.
\end{align*}
\[
 H_3=B\cap C,\qquad
 P(H_3)=P(B)P(C\mid B)=\frac{30}{40}\frac{29}{39}
 =\frac{870}{40\cdot39}.
\]

\parttitle{1. Вероятность сдачи экзамена}
\[
 P(A\mid H_1)=0,\qquad P(A\mid H_2)=\frac{29}{38},\qquad
 P(A\mid H_3)=1.
\]
При $H_2$ из оставшихся $38$ вопросов студент знает $29$;
это и даёт условную вероятность ответа на дополнительный вопрос.
По формуле полной вероятности
\begin{align*}
 P(A)&=P(H_1)P(A\mid H_1)+P(H_2)P(A\mid H_2)+P(H_3)P(A\mid H_3)\\
 &=\frac{90}{40\cdot39}\cdot0
 +\frac{600}{40\cdot39}\frac{29}{38}
 +\frac{870}{40\cdot39}\cdot1\\
 &=\frac{841}{988}\approx0{,}851.
\end{align*}
\editorial{В рукописи результат округлён до $0{,}852$.
Для записанной суммы более точное значение $0{,}851214\ldots$;
округление до трёх знаков исправлено.}

\sourcepage{rk2-additions/photo_2026-11-07_23-15-51_1002843.jpg}
\parttitle{2. Ровно один известный вопрос при условии сдачи экзамена}
По формуле Байеса
\[
 P(H_2\mid A)=\frac{P(H_2)P(A\mid H_2)}{P(A)}
 =\frac{\dfrac{600}{40\cdot39}\dfrac{29}{38}}{841/988}
 =\frac{10}{29}\approx0{,}345.
\]
\end{document}
